\appendix

\section{Proof Details}

\subsection{Inner derivative and image length}
\label{app:inner-certificate}

\begin{lemma}[Inner-sweep derivative and Lipschitz control]
\label{lem:inner-derivative}
Under Assumptions~\ref{assump:indexed-regime} and
\ref{assump:compact-cube-support}, if $d\geq1$, $R\geq1$, and
$\alpha_{1:d-1}\in[-R,R]^{d-1}$, then for every $\theta\in[-1,1]$,
\[
\partial_\theta s_0(\theta;\alpha_{1:d-1})
=-d\theta^{d-1}-\sum_{j=1}^{d-1}j\alpha_j\theta^{j-1}
\]
and
\[
\bigl|\partial_\theta s_0(\theta;\alpha_{1:d-1})\bigr|
\leq d+\frac{Rd(d-1)}2=B_0(d,R).
\]
Consequently, for all $x,z\in[-1,1]$,
\[
|s_0(x;\alpha_{1:d-1})-s_0(z;\alpha_{1:d-1})|
\leq B_0(d,R)|x-z|.
\]
\end{lemma}

\begin{proof}
The definition
\[
s_0(\theta;\alpha_{1:d-1})
=-\theta^d-\sum_{j=1}^{d-1}\alpha_j\theta^j
\]
is polynomial in $\theta$, so termwise differentiation gives the asserted
formula on all of $\R$.  First suppose $d\geq2$.  For
$\theta\in[-1,1]$, Assumption~\ref{assump:compact-cube-support} and the
triangle inequality give
\[
\begin{aligned}
\bigl|\partial_\theta s_0(\theta;\alpha_{1:d-1})\bigr|
&\leq d|\theta|^{d-1}
 +\sum_{j=1}^{d-1}j|\alpha_j||\theta|^{j-1}\\
&\leq d+R\sum_{j=1}^{d-1}j\\
&=d+\frac{Rd(d-1)}2
=B_0(d,R).
\end{aligned}
\]
No term has been absorbed.  For distinct $x,z\in[-1,1]$, the closed segment
between them lies in $[-1,1]$; the mean-value theorem therefore yields the
Lipschitz inequality.  It is immediate when $x=z$, and the argument includes
both endpoints $-1$ and $1$.

When $d=1$, the coefficient tuple and every displayed sum are empty.  Without
using a $0^0$ convention,
\[
s_0(\theta)=-\theta,\qquad s_0'(\theta)=-1,
\qquad B_0(1,R)=1,
\]
so the derivative and Lipschitz conclusions remain exact.  At the radius
boundary $R=1$, the same computation uses $|\alpha_j|\leq1$ and gives
$B_0(d,1)=d+d(d-1)/2$; no strict radius inequality or limiting argument is
needed.
\end{proof}

\begin{lemma}[Inner interval image length]
\label{lem:inner-image-length}
Under Assumptions~\ref{assump:indexed-regime} and
\ref{assump:compact-cube-support} and Lemma~\ref{lem:inner-derivative}, if
$d\geq1$, $R\geq1$, $\alpha_{1:d-1}\in[-R,R]^{d-1}$, and
$J\subseteq[-1,1]$ is an interval with any endpoint convention, then
$s_0(J;\alpha_{1:d-1})$ is an interval or is empty, and
\[
\lambda\bigl(s_0(J;\alpha_{1:d-1})\bigr)
\leq B_0(d,R)|J|.
\]
This includes empty and singleton $J$, intervals containing either boundary
point, and $d=1$.
\end{lemma}

\begin{proof}
If $J=\varnothing$, its image is empty and both lengths are zero.  If
$J=\{\theta_0\}$, its image is a singleton and again both lengths are zero.
Suppose now that $J$ has at least two points.  Every real interval is
connected, regardless of endpoint inclusion.  Since $s_0$ is continuous,
its image is connected and hence is an interval in $\R$.  It is bounded
because Lemma~\ref{lem:inner-derivative} makes $s_0$ Lipschitz on the bounded
set $[-1,1]$.

For arbitrary $u,v$ in the image, choose $x,z\in J$ with
$u=s_0(x;\alpha_{1:d-1})$ and $v=s_0(z;\alpha_{1:d-1})$.  Then
\[
|u-v|\leq B_0(d,R)|x-z|
\leq B_0(d,R)\diam(J).
\]
Taking the supremum over $u,v$ gives
\[
\diam\bigl(s_0(J;\alpha_{1:d-1})\bigr)
\leq B_0(d,R)\diam(J).
\]
The Lebesgue length of any nonempty bounded real interval equals its diameter,
independently of whether either endpoint is included.  Hence
\[
\begin{aligned}
\lambda\bigl(s_0(J;\alpha_{1:d-1})\bigr)
&=\diam\bigl(s_0(J;\alpha_{1:d-1})\bigr)\\
&\leq B_0(d,R)\diam(J)
=B_0(d,R)|J|.
\end{aligned}
\]
This proves the complete inner deterministic certificate.  No monotonicity,
injectivity, transversality, or simple-root condition is used.  In degree one
the image is exactly $-J$, so equality holds with $B_0(1,R)=1$.
\end{proof}

\subsection{Positive and negative outer certificates}
\label{app:outer-certificates}

\begin{proposition}[Positive outer pivot and image certificate]
\label{prop:positive-outer-certificate}
Under Assumptions~\ref{assump:indexed-regime} and
\ref{assump:compact-cube-support}, fix $\alpha\in[-R,R]^d$.  If
$J\subseteq(1,\infty)$ is an interval with any endpoint convention, including
an empty or singleton interval, then for every $\theta\in J$,
\[
\phi_\alpha(\theta)=0
\quad\Longleftrightarrow\quad
\alpha_{d-1}=s_\infty(\theta;\alpha_{0:d-2})
=-\theta-\sum_{j=0}^{d-2}\alpha_j\theta^{j-d+1}.
\]
Moreover,
\[
\partial_\theta s_\infty(\theta;\alpha_{0:d-2})
=-1+\sum_{j=0}^{d-2}(d-1-j)\alpha_j\theta^{j-d},
\]
\[
\bigl|\partial_\theta s_\infty(\theta;\alpha_{0:d-2})\bigr|
\leq B_\infty(d,R),\qquad B_\infty(d,R)=1+\frac{Rd(d-1)}2,
\]
and
\[
\lambda\bigl(s_\infty(J;\alpha_{0:d-2})\bigr)
\leq B_\infty(d,R)|J|.
\]
The conclusions remain valid as an endpoint approaches $1$ from above, and
for $d=1$ they reduce to $s_\infty=-\theta$ and exact image length $|J|$.
\end{proposition}

\begin{proof}
For $\theta>1$, the divisor $\theta^{d-1}$ is nonzero.  Separating the
coefficient of degree $d-1$ gives the reversible chain
\[
\begin{aligned}
\phi_\alpha(\theta)=0
&\Longleftrightarrow \frac{\phi_\alpha(\theta)}{\theta^{d-1}}=0\\
&\Longleftrightarrow
\theta+\sum_{j=0}^{d-2}\alpha_j\theta^{j-d+1}
+\alpha_{d-1}=0\\
&\Longleftrightarrow
\alpha_{d-1}=-\theta-\sum_{j=0}^{d-2}\alpha_j\theta^{j-d+1}.
\end{aligned}
\]
Termwise differentiation yields
\[
\begin{aligned}
\partial_\theta s_\infty(\theta;\alpha_{0:d-2})
&=-1-\sum_{j=0}^{d-2}(j-d+1)\alpha_j\theta^{j-d}\\
&=-1+\sum_{j=0}^{d-2}(d-1-j)\alpha_j\theta^{j-d}.
\end{aligned}
\]
For $d\geq2$, every exponent in the last sum satisfies $j-d\leq-2$ and
$|\theta^{j-d}|\leq1$.  Therefore
\[
\begin{aligned}
|\partial_\theta s_\infty(\theta;\alpha_{0:d-2})|
&\leq1+\sum_{j=0}^{d-2}(d-1-j)|\alpha_j||\theta^{j-d}|\\
&\leq1+R\sum_{j=0}^{d-2}(d-1-j)\\
&=1+R\sum_{k=1}^{d-1}k
=1+\frac{Rd(d-1)}2=B_\infty(d,R).
\end{aligned}
\]
The final identity also covers $d=1$, when both sums are empty.

If $|J|=\infty$, the image-length inequality holds in the extended
nonnegative reals.  If $J$ is empty or a singleton, both sides have length
zero.  Otherwise, every segment joining two points of $J$ remains in
$(1,\infty)$.  The mean-value theorem gives the
$B_\infty(d,R)$-Lipschitz bound there.  The continuous image of the connected
set $J$ is an interval, and the length-diameter identity for real intervals
gives
\[
\begin{aligned}
\lambda\bigl(s_\infty(J;\alpha_{0:d-2})\bigr)
&=\diam\bigl(s_\infty(J;\alpha_{0:d-2})\bigr)\\
&\leq B_\infty(d,R)\diam(J)
=B_\infty(d,R)|J|.
\end{aligned}
\]
This is unchanged by inclusion or exclusion of a finite endpoint.

The one-sided limits
\[
\begin{aligned}
\lim_{\theta\downarrow1}s_\infty(\theta;\alpha_{0:d-2})
&=-1-\sum_{j=0}^{d-2}\alpha_j,\\
\lim_{\theta\downarrow1}\partial_\theta s_\infty(\theta;\alpha_{0:d-2})
&=-1+\sum_{j=0}^{d-2}(d-1-j)\alpha_j
\end{aligned}
\]
are finite and satisfy the same derivative bound.  Thus approaching the
excluded endpoint adds no length term; the point $1$ itself belongs to the
inner chart.  When $d=1$, division is by $\theta^0=1$, all sums are empty,
$\phi_\alpha(\theta)=\theta+\alpha_0$, and all stated identities are exact.
\end{proof}

\begin{proposition}[Negative outer pivot and image certificate]
\label{prop:negative-outer-certificate}
Under Assumptions~\ref{assump:indexed-regime} and
\ref{assump:compact-cube-support}, fix $\alpha\in[-R,R]^d$.  If
$J\subseteq(-\infty,-1)$ is an interval with any endpoint convention,
including an empty or singleton interval, then for every $\theta\in J$,
\[
\phi_\alpha(\theta)=0
\quad\Longleftrightarrow\quad
\alpha_{d-1}=s_\infty(\theta;\alpha_{0:d-2})
=-\theta-\sum_{j=0}^{d-2}\alpha_j\theta^{j-d+1}.
\]
The exact derivative is
\[
\partial_\theta s_\infty(\theta;\alpha_{0:d-2})
=-1+\sum_{j=0}^{d-2}(d-1-j)\alpha_j\theta^{j-d},
\]
and it satisfies
\[
|\partial_\theta s_\infty(\theta;\alpha_{0:d-2})|
\leq B_\infty(d,R),
\qquad
\lambda\bigl(s_\infty(J;\alpha_{0:d-2})\bigr)
\leq B_\infty(d,R)|J|.
\]
These conclusions preserve the negative-$\theta$ parity signs, remain valid
as an endpoint approaches $-1$ from below, and specialize exactly at $d=1$.
\end{proposition}

\begin{proof}
For $\theta<-1$, the number $\theta^{d-1}$ is nonzero.  It may be negative,
but it divides an equality rather than an inequality, so
\[
\begin{aligned}
\phi_\alpha(\theta)=0
&\Longleftrightarrow \frac{\phi_\alpha(\theta)}{\theta^{d-1}}=0\\
&\Longleftrightarrow
\theta+\sum_{j=0}^{d-2}\alpha_j\theta^{j-d+1}
+\alpha_{d-1}=0\\
&\Longleftrightarrow
\alpha_{d-1}=-\theta-\sum_{j=0}^{d-2}\alpha_j\theta^{j-d+1}.
\end{aligned}
\]
Multiplication by the same nonzero divisor proves the converse.  Direct
differentiation gives the derivative displayed in the proposition.

To retain every parity sign, orient this connected chart by $u=-\theta>1$
and put
\[
g(u):=s_\infty(-u;\alpha_{0:d-2})
=u-\sum_{j=0}^{d-2}\alpha_j(-u)^{j-d+1}.
\]
The chain rule gives
\[
\begin{aligned}
g'(u)
&=-\partial_\theta s_\infty(-u;\alpha_{0:d-2})\\
&=1-\sum_{j=0}^{d-2}(d-1-j)\alpha_j(-u)^{j-d}\\
&=1-\sum_{j=0}^{d-2}(d-1-j)\alpha_j
(-1)^{d-j}u^{j-d}.
\end{aligned}
\]
Here $(-u)^{j-d}=(-1)^{j-d}u^{j-d}=(-1)^{d-j}u^{j-d}$ because the exponents
are integers.  Thus orientation reversal gives exactly
$g'(u)=-s_\infty'(-u)$ and assigns no favorable sign to a coefficient term.

For $d\geq2$, $j-d\leq-2$ and
$|\theta^{j-d}|=|\theta|^{-(d-j)}\leq1$.  Consequently,
\[
\begin{aligned}
|\partial_\theta s_\infty(\theta;\alpha_{0:d-2})|
&\leq1+\sum_{j=0}^{d-2}(d-1-j)|\alpha_j||\theta^{j-d}|\\
&\leq1+R\sum_{j=0}^{d-2}(d-1-j)\\
&=1+R\sum_{k=1}^{d-1}k
=1+\frac{Rd(d-1)}2=B_\infty(d,R).
\end{aligned}
\]
The same bound holds for $|g'(u)|$, and the triangular-sum identity includes
the empty $d=1$ sum.

If $|J|=\infty$, the length conclusion is immediate in the extended reals.
Empty and singleton intervals give $0\leq0$.  Otherwise, every segment
joining points of $J$ stays inside $(-\infty,-1)$ and crosses neither $-1$
nor zero.  The mean-value theorem, connectedness, and the length-diameter
identity therefore give
\[
\begin{aligned}
\lambda\bigl(s_\infty(J;\alpha_{0:d-2})\bigr)
&=\diam\bigl(s_\infty(J;\alpha_{0:d-2})\bigr)\\
&\leq B_\infty(d,R)\diam(J)
=B_\infty(d,R)|J|.
\end{aligned}
\]

Finally, the relevant one-sided limits are
\[
\begin{aligned}
\lim_{\theta\uparrow-1}s_\infty(\theta;\alpha_{0:d-2})
&=1-\sum_{j=0}^{d-2}\alpha_j(-1)^{j-d+1},\\
\lim_{\theta\uparrow-1}\partial_\theta s_\infty(\theta;\alpha_{0:d-2})
&=-1+\sum_{j=0}^{d-2}(d-1-j)\alpha_j(-1)^{j-d},\\
\lim_{u\downarrow1}g'(u)
&=1-\sum_{j=0}^{d-2}(d-1-j)\alpha_j(-1)^{j-d}.
\end{aligned}
\]
The last two quantities are negatives and each has absolute value at most
$B_\infty(d,R)$.  Thus approaching $-1$ introduces neither a blow-up nor an
endpoint correction; $-1$ belongs to the inner chart.  In degree one all
sums are empty, division is by $1$, and $s_\infty=-\theta$ on this chart as
well.  The two outer propositions have treated the connected components
separately and nowhere join points across zero.
\end{proof}

\subsection{Measurable caps, Borel sections, and matching disintegration}
\label{app:disintegration}

\begin{lemma}[Measurable endpoint-kernel caps]
\label{lem:kernel-caps}
Under Assumptions~\ref{assump:indexed-regime},
\ref{assump:compact-cube-support}, and
\ref{assump:mean-endpoint-conditional-caps}, let
$\mu\in\mathcal D_{d,R,\eta}$.  For $d\geq2$, define
\[
(X_0,Y_0):=(\alpha_0,\alpha_{1:d-1}),\qquad
(X_\infty,Y_\infty):=(\alpha_{d-1},\alpha_{0:d-2}),
\]
let $\nu_i$ be the law of $Y_i$, and fix a regular conditional kernel
$Q^i_y=\mathcal L(X_i\mid Y_i=y)$ for $i\in\{0,\infty\}$.  For $d=1$,
take both pairs to be $(\alpha_0,\varnothing)$ and use the same unconditional
kernel on $\R^0=\{\varnothing\}$.

There exist extended nonnegative measurable functions $\widehat K_i$ and
$\nu_i$-full measurable sets $G_i$ such that every $y\in G_i$ is
support-compatible and
\[
Q^i_y\ll\lambda,\qquad Q^i_y([-R,R])=1,\qquad
\widehat K_i(y)=\left\|\frac{\dd Q^i_y}{\dd\lambda}\right\|_{L^\infty(\R)}
<\infty.
\]
On the same set, simultaneously for every Borel $A\subseteq\R$,
\begin{equation}
Q^i_y(A)\leq\widehat K_i(y)\lambda(A).
\label{eq:kernel-domination}
\end{equation}
Furthermore,
\begin{equation}
\widehat K_0(Y_0)=K_0^\mu,\qquad
\widehat K_\infty(Y_\infty)=K_\infty^\mu
\quad\mu\text{-a.s.},
\label{eq:cap-identities}
\end{equation}
and
\begin{equation}
\mathbb E_\mu\widehat K_0(Y_0)\leq\bar\kappa_0,
\qquad
\mathbb E_\mu\widehat K_\infty(Y_\infty)\leq\bar\kappa_\infty.
\label{eq:mean-cap-bounds}
\end{equation}
The representatives are unchanged almost surely if a conditional-kernel or
density version is changed.  In degree one the two kernels and caps coincide.
\end{lemma}

\begin{proof}
Fix one index $i$, retaining its own conditioning marginal $\nu_i$ and
kernel $Q^i$ throughout.  The argument will then be applied independently to
$i=0$ and $i=\infty$.

By Assumption~\ref{assump:mean-endpoint-conditional-caps}, $Q^i_y\ll\lambda$
outside a $\nu_i$-null set.  On the product of the conditioning and pivot
spaces define
\[
\gamma_i(C):=\int Q^i_y(C_y)\,\nu_i(\dd y),
\qquad C_y:=\{x:(y,x)\in C\}.
\]
If $(\nu_i\otimes\lambda)(C)=0$, Tonelli's theorem gives
$\lambda(C_y)=0$ for $\nu_i$-almost every $y$.  Conditional absolute
continuity then gives $Q^i_y(C_y)=0$, and hence
$\gamma_i\ll\nu_i\otimes\lambda$.  The Radon--Nikodym theorem supplies a
nonnegative jointly measurable $q_i(y,x)$ such that
\[
\gamma_i(\dd y,\dd x)=q_i(y,x)\nu_i(\dd y)\lambda(\dd x).
\]
For every member $B$ of a countable $\pi$-system generating
$\mathcal B(\R)$, comparison on all rectangles $D\times B$ gives
\[
Q^i_y(B)=\int_Bq_i(y,x)\dd x
\quad\text{for $\nu_i$-almost every $y$}.
\]
Intersecting the countably many full-measure sets and using uniqueness of
finite measures extends this identity to every Borel $B$ on one full-measure
set.  Thus $q_i(y,\cdot)$ is a density of $Q^i_y$ there.

Let
\[
\mathscr J_{\mathbb Q}:=\{(a,b):a,b\in\mathbb Q, a<b\},
\qquad
\widehat K_i(y):=\sup_{J\in\mathscr J_{\mathbb Q}}
\frac{Q^i_y(J)}{|J|}.
\]
For each fixed $J$, the kernel property makes $y\mapsto Q^i_y(J)$
measurable.  The family is countable, so $\widehat K_i$ is measurable on the
whole conditioning space, including fibers on which no density was chosen.

Fix a density fiber, write its density as $q$, and set
$m=\|q\|_{L^\infty(\R)}$, allowing $m=+\infty$.  For every rational interval,
\[
\frac{Q^i_y(J)}{|J|}=\frac1{|J|}\int_Jq(x)\dd x\leq m,
\]
so $\widehat K_i(y)\leq m$.  Conversely, take a finite $a<m$.  The set
$\{x:q(x)>a\}$ has positive measure and therefore contains a Lebesgue point
of $q$.  Some bounded open interval $L$ around that point satisfies
\[
\frac1{|L|}\int_Lq(t)\dd t>a.
\]
Approximating both endpoints of $L$ by rationals preserves this strict
inequality, by absolute continuity of the integral of the $L^1$ density and
convergence of interval lengths.  Thus $\widehat K_i(y)>a$.  Letting
$a\uparrow m$, or taking arbitrarily large finite $a$ when $m=+\infty$, gives
\begin{equation}
\widehat K_i(y)=\|q\|_{L^\infty(\R)}
\label{eq:rational-cap-identity}
\end{equation}
on every density fiber.

Two densities of the same fiber law agree $\lambda$-almost everywhere, so
their essential suprema agree.  If $Q^i$ and $\widetilde Q^i$ are two regular
conditional kernels, conditional-expectation uniqueness gives
$Q^i_y(J)=\widetilde Q^i_y(J)$ for $\nu_i$-almost every $y$ for each
$J\in\mathscr J_{\mathbb Q}$.  Countability yields one common full-measure
set, and the rational-interval formula gives the same cap there.  This proves
both kinds of version independence.

On a density fiber with finite cap,
$q(x)\leq\widehat K_i(y)$ for $\lambda$-almost every $x$.  Consequently,
for every Borel $A$ on that same fiber, with no $A$-dependent exceptional
set,
\[
Q^i_y(A)=\int_Aq(x)\dd x
\leq\widehat K_i(y)\lambda(A).
\]
Equation~\eqref{eq:rational-cap-identity} identifies the random cap with the
one in the definition of $\mathcal D_{d,R,\eta}$, proving
\eqref{eq:cap-identities}.  The primitive mean inequalities give
\eqref{eq:mean-cap-bounds} without changing constants.  Their finite right
sides imply that both caps are finite almost surely.

It remains to put conditional support on the same full-measure set.  By
Assumption~\ref{assump:compact-cube-support},
\[
0=\mu(X_i\notin[-R,R])
=\int Q^i_y(\R\setminus[-R,R])\nu_i(\dd y).
\]
The nonnegative integrand vanishes almost everywhere, and $Y_i$ lies in its
corresponding $(d-1)$-dimensional cube almost surely.  Intersecting the
full-measure sets for density, finite cap, and support produces $G_i$.
Applying the construction separately to $0$ and $\infty$ proves the result
for $d\geq2$.

For $d=1$, the conditioning marginal is the unit mass on $\R^0$.  The kernel
at its sole point is the unconditional law of $\alpha_0$, and the same kernel
is used for both endpoint labels.  The rational-interval formula therefore
produces one cap, and all preceding identities give both required
inequalities for it.
\end{proof}

\begin{lemma}[Borel root events and exact chart sections]
\label{lem:borel-sections}
Under Assumptions~\ref{assump:compact-parameter-domain},
\ref{assump:indexed-regime}, and \ref{assump:compact-cube-support}, let
$J\subseteq\Theta$ be an interval with any endpoint convention, possibly
empty or a singleton.  Then $H_{d,J}$ is Borel.  For $d\geq2$, define
\[
T_0(x,y):=(x,y_1,\ldots,y_{d-1}),\qquad
T_\infty(x,y):=(y_0,\ldots,y_{d-2},x),
\]
and $E^i_J:=T_i^{-1}(H_{d,J})$.  For $d=1$, define both maps by
$T_i(x,\varnothing)=(x)$.  The sets $E^i_J$ and all their pivot sections
$(E^i_J)_y:=\{x:(x,y)\in E^i_J\}$ are Borel.  If $J\subseteq[-1,1]$, then
\begin{equation}
(E^0_J)_y=
\begin{cases}
[-R,R]\cap s_0(J;y),&y\in[-R,R]^{d-1},\\
\varnothing,&y\notin[-R,R]^{d-1},
\end{cases}
\label{eq:inner-section}
\end{equation}
and if $J\subseteq(1,\infty)$ or $J\subseteq(-\infty,-1)$, then
\begin{equation}
(E^\infty_J)_y=
\begin{cases}
[-R,R]\cap s_\infty(J;y),&y\in[-R,R]^{d-1},\\
\varnothing,&y\notin[-R,R]^{d-1}.
\end{cases}
\label{eq:outer-section}
\end{equation}
For $d=1$, both right sides mean $[-R,R]\cap(-J)$.
\end{lemma}

\begin{proof}
Assumption~\ref{assump:compact-parameter-domain} makes $J$ bounded.  The
empty interval is a countable union of empty compact sets, and a singleton is
compact.  Otherwise write $a=\inf J<b=\sup J$.  For all sufficiently large
$n$, set
\[
a_n:=\begin{cases}a,&a\in J,\\a+1/n,&a\notin J,\end{cases}
\qquad
b_n:=\begin{cases}b,&b\in J,\\b-1/n,&b\notin J,\end{cases}
\]
and let $K_n=[a_n,b_n]$ when $a_n\leq b_n$, taking $K_n=\varnothing$
otherwise.  After reindexing, the compact sets $K_n$ increase to $J$.  This
construction preserves open, closed, and half-open endpoint conventions.

For every $n$, define
\[
Z_n:=\{(\alpha,\theta)\in[-R,R]^d\times K_n:
\phi_\alpha(\theta)=0\}.
\]
The polynomial evaluation map is continuous, so $Z_n$ is closed in a compact
space and hence compact.  Its coefficient projection is compact and Borel,
and exact existential quantification gives
\[
H_{d,J}=\bigcup_{n=1}^\infty\operatorname{proj}_\alpha Z_n.
\]
Thus $H_{d,J}$ is Borel.  The proof projects compact sets before taking a
countable union; it does not rely on projections of arbitrary Borel sets.

The reconstruction maps are continuous, so each $E^i_J$ is Borel.  Sections
of product Borel sets are Borel: rectangles have Borel sections, and this
property is closed under complements and countable unions.

For an inner interval and support-compatible $y$,
\[
\phi_{T_0(x,y)}(\theta)=0
\Longleftrightarrow
x=-\theta^d-\sum_{j=1}^{d-1}y_j\theta^j=s_0(\theta;y).
\]
The coefficient cube then gives \eqref{eq:inner-section}; an incompatible
$y$ gives an empty section.  On either outer chart, $\theta\neq0$, and
\[
\begin{aligned}
\phi_{T_\infty(x,y)}(\theta)=0
&\Longleftrightarrow
\theta+\sum_{j=0}^{d-2}y_j\theta^{j-d+1}+x=0\\
&\Longleftrightarrow x=s_\infty(\theta;y),
\end{aligned}
\]
where division is only by the nonzero $\theta^{d-1}$.  This proves
\eqref{eq:outer-section} on both signs.  In degree one both sums are empty
and both reconstructions are $(x)$.

Only the zero-set identity is used, so tangencies and multiple roots need no
extra case.  Empty and singleton intervals and every endpoint convention
were included in the compact exhaustion.  In particular, $\pm1$ occur only
in the inner chart when they belong to $J$.
\end{proof}

\begin{proposition}[Matching-kernel Borel disintegration]
\label{prop:matching-disintegration}
Under Assumptions~\ref{assump:compact-parameter-domain}--
\ref{assump:mean-endpoint-conditional-caps} and
Lemmas~\ref{lem:kernel-caps} and \ref{lem:borel-sections}, use the pairs,
kernels, marginals, full-measure sets, and reconstruction events defined
there.  For every inner-chart interval $J\subseteq\Theta\cap[-1,1]$,
\begin{equation}
\mu(H_{d,J})=\int Q^0_y((E^0_J)_y)\nu_0(\dd y),
\label{eq:inner-disintegration}
\end{equation}
the integrand is measurable, and for every $y\in G_0$,
\begin{equation}
Q^0_y((E^0_J)_y)=Q^0_y(s_0(J;y))
\leq\widehat K_0(y)\lambda(s_0(J;y)).
\label{eq:inner-fiber-bound}
\end{equation}
For every interval $J$ in either outer chart,
\begin{equation}
\mu(H_{d,J})=\int Q^\infty_y((E^\infty_J)_y)\nu_\infty(\dd y),
\label{eq:outer-disintegration}
\end{equation}
the integrand is measurable, and for every $y\in G_\infty$,
\begin{equation}
Q^\infty_y((E^\infty_J)_y)=Q^\infty_y(s_\infty(J;y))
\leq\widehat K_\infty(y)\lambda(s_\infty(J;y)).
\label{eq:outer-fiber-bound}
\end{equation}
For $d\geq2$ the two displays use their distinct conditioning kernels; for
$d=1$ they use the single unconditional kernel of $\alpha_0$.
\end{proposition}

\begin{proof}
We first prove the required parameterized measurability.  For a probability
kernel $Q_y$ from a measurable space $S$ to $\R$, let $\mathcal M$ be the
class of Borel sets $E\subseteq\R\times S$ for which
$y\mapsto Q_y(E_y)$ is measurable.  Rectangles belong to $\mathcal M$
because
\[
Q_y((A\times B)_y)=\1_B(y)Q_y(A).
\]
The class is closed under complements and countable disjoint unions, since
\[
Q_y((E^c)_y)=1-Q_y(E_y),\qquad
Q_y\left(\left(\bigcup_nE_n\right)_y\right)=\sum_nQ_y((E_n)_y).
\]
Thus $\mathcal M$ is a $\lambda$-system containing the $\pi$-system of
measurable rectangles.  The $\pi$--$\lambda$ theorem proves the assertion
for every product-Borel $E$.

If $Q_y$ is a regular conditional law of $X$ given $Y$, the identity
\[
\mathbb P((X,Y)\in A\times B)=\int_BQ_y(A)\nu(\dd y)
\]
holds on rectangles.  The sets on which the corresponding integral identity
holds form another $\lambda$-system: complements follow by subtraction from
one, while disjoint unions follow by countable additivity and monotone
convergence.  A second $\pi$--$\lambda$ argument gives, for every Borel
$E\subseteq\R\times S$,
\begin{equation}
\mathbb P((X,Y)\in E)=\int Q_y(E_y)\nu(\dd y).
\label{eq:general-disintegration}
\end{equation}

Apply \eqref{eq:general-disintegration} to $(X_0,Y_0)$ and the Borel set
$E^0_J$.  Since $T_0(X_0,Y_0)=\alpha$, this proves
\eqref{eq:inner-disintegration}.  For $y\in G_0$,
Lemma~\ref{lem:borel-sections} gives
$(E^0_J)_y=[-R,R]\cap s_0(J;y)$, while Lemma~\ref{lem:kernel-caps} gives
$Q^0_y([-R,R])=1$.  Intersecting with the cube therefore does not change
conditional probability.  The continuous image of an interval is an
interval and hence Borel, so \eqref{eq:kernel-domination} applied to
$s_0(J;y)$ proves \eqref{eq:inner-fiber-bound}.

Apply the same general identity separately to $(X_\infty,Y_\infty)$ and
$E^\infty_J$.  The outer section formula, support of $Q^\infty_y$, and the
same Borel-set domination prove \eqref{eq:outer-disintegration} and
\eqref{eq:outer-fiber-bound}.  No inner kernel enters this calculation, and
no outer kernel enters the inner calculation.

When $d=1$, both coordinate pairs are $(\alpha_0,\varnothing)$, both
reconstruction maps coincide, and the conditioning marginal is a unit mass.
The kernel identification in Lemma~\ref{lem:kernel-caps} then makes all four
displays the correct unconditional identities.  Empty chart pieces give
empty images, and singleton pieces give singleton images; conditional
absolute continuity assigns the latter probability zero.  This completes
the measurable conditional interface used below.
\end{proof}

\subsection{Inner, positive, and negative chart probability controls}
\label{app:chart-probabilities}

\begin{proposition}[Inner chart probability control]
\label{prop:inner-chart-probability}
Under Assumptions~\ref{assump:indexed-regime},
\ref{assump:compact-cube-support}, and
\ref{assump:mean-endpoint-conditional-caps}, and under
Lemmas~\ref{lem:inner-image-length} and \ref{lem:kernel-caps} and
Proposition~\ref{prop:matching-disintegration}, fix
$\mu\in\mathcal D_{d,R,\eta}$ and $I\in\mathcal I(\Theta)$.  With
$I_0=I\cap[-1,1]$,
\[
\mu(H_{d,I_0})\leq\bar\kappa_0B_0(d,R)|I_0|.
\]
If $I_0$ is empty or a singleton, then $\mu(H_{d,I_0})=0$.
\end{proposition}

\begin{proof}
Proposition~\ref{prop:matching-disintegration} gives a measurable integrand,
and the full-measure set $G_0$ from Lemma~\ref{lem:kernel-caps} permits us to
write, without using a cap on a null fiber,
\[
\begin{aligned}
\mu(H_{d,I_0})
&=\int_{G_0}Q^0_y((E^0_{I_0})_y)\nu_0(\dd y)\\
&=\int_{G_0}Q^0_y(s_0(I_0;y))\nu_0(\dd y).
\end{aligned}
\]
For every $y\in G_0$, Lemma~\ref{lem:kernel-caps} supplies finite
$\widehat K_0(y)$ and simultaneous Borel-set domination, while
Lemma~\ref{lem:inner-image-length} supplies
\[
\lambda(s_0(I_0;y))\leq B_0(d,R)|I_0|.
\]
Consequently,
\[
\begin{aligned}
\mu(H_{d,I_0})
&\leq\int_{G_0}\widehat K_0(y)\lambda(s_0(I_0;y))\nu_0(\dd y)\\
&\leq B_0(d,R)|I_0|\int_{G_0}\widehat K_0(y)\nu_0(\dd y)\\
&\leq B_0(d,R)|I_0|\mathbb E_\mu\widehat K_0(Y_0)\\
&\leq\bar\kappa_0B_0(d,R)|I_0|.
\end{aligned}
\]
If $I_0=\varnothing$, the event is empty.  If $I_0$ is a singleton, its
image is a singleton and has Lebesgue measure zero; on $G_0$ the finite cap
then gives conditional probability zero.  Thus the degenerate conclusion
does not involve a zero-times-infinity product, even when
$\bar\kappa_0=0$.
\end{proof}

\begin{proposition}[Positive outer chart probability control]
\label{prop:positive-chart-probability}
Under Assumptions~\ref{assump:indexed-regime},
\ref{assump:compact-cube-support}, and
\ref{assump:mean-endpoint-conditional-caps}, and under
Propositions~\ref{prop:positive-outer-certificate},
\ref{prop:matching-disintegration}, and Lemma~\ref{lem:kernel-caps}, fix
$\mu\in\mathcal D_{d,R,\eta}$ and $I\in\mathcal I(\Theta)$.  With
$I_+=I\cap(1,\infty)$,
\[
\mu(H_{d,I_+})\leq\bar\kappa_\infty B_\infty(d,R)|I_+|.
\]
If $I_+$ is empty or a singleton, then $\mu(H_{d,I_+})=0$.
\end{proposition}

\begin{proof}
The positive-chart instance of Proposition~\ref{prop:matching-disintegration}
and the full-measure property of $G_\infty$ give
\[
\begin{aligned}
\mu(H_{d,I_+})
&=\int_{G_\infty}Q^\infty_y((E^\infty_{I_+})_y)\nu_\infty(\dd y)\\
&=\int_{G_\infty}Q^\infty_y(s_\infty(I_+;y))\nu_\infty(\dd y).
\end{aligned}
\]
For $d\geq2$, $y\in[-R,R]^{d-1}$ on $G_\infty$; appending any pivot value
in the coefficient cube makes the deterministic outer certificate applicable.
For $d=1$ the empty-tuple specialization applies directly.  Thus
Proposition~\ref{prop:positive-outer-certificate} gives
\[
\lambda(s_\infty(I_+;y))\leq B_\infty(d,R)|I_+|.
\]
Using the finite cap and then its mean bound yields
\[
\begin{aligned}
\mu(H_{d,I_+})
&\leq\int_{G_\infty}\widehat K_\infty(y)
       \lambda(s_\infty(I_+;y))\nu_\infty(\dd y)\\
&\leq B_\infty(d,R)|I_+|
       \int_{G_\infty}\widehat K_\infty(y)\nu_\infty(\dd y)\\
&\leq B_\infty(d,R)|I_+|
       \mathbb E_\mu\widehat K_\infty(Y_\infty)\\
&\leq\bar\kappa_\infty B_\infty(d,R)|I_+|.
\end{aligned}
\]
For an empty piece the event is empty.  For a singleton, the image is a
singleton and the finite-cap domination on $G_\infty$ gives zero conditional
probability.  No conditioning-null fiber is used.
\end{proof}

\begin{proposition}[Negative outer chart probability control]
\label{prop:negative-chart-probability}
Under Assumptions~\ref{assump:indexed-regime},
\ref{assump:compact-cube-support}, and
\ref{assump:mean-endpoint-conditional-caps}, and under
Propositions~\ref{prop:negative-outer-certificate},
\ref{prop:matching-disintegration}, and Lemma~\ref{lem:kernel-caps}, fix
$\mu\in\mathcal D_{d,R,\eta}$ and $I\in\mathcal I(\Theta)$.  With
$I_-=I\cap(-\infty,-1)$,
\[
\mu(H_{d,I_-})\leq\bar\kappa_\infty B_\infty(d,R)|I_-|.
\]
If $I_-$ is empty or a singleton, then $\mu(H_{d,I_-})=0$.
\end{proposition}

\begin{proof}
The negative-chart instance of Proposition~\ref{prop:matching-disintegration}
gives, after restriction to the finite-cap set,
\[
\begin{aligned}
\mu(H_{d,I_-})
&=\int_{G_\infty}Q^\infty_y((E^\infty_{I_-})_y)\nu_\infty(\dd y)\\
&=\int_{G_\infty}Q^\infty_y(s_\infty(I_-;y))\nu_\infty(\dd y).
\end{aligned}
\]
Appending a pivot value to a support-compatible $y$ and applying the
specifically negative certificate gives
\[
\lambda(s_\infty(I_-;y))\leq B_\infty(d,R)|I_-|,
\]
including the empty-tuple case $d=1$.  Therefore
\[
\begin{aligned}
\mu(H_{d,I_-})
&\leq\int_{G_\infty}\widehat K_\infty(y)
       \lambda(s_\infty(I_-;y))\nu_\infty(\dd y)\\
&\leq B_\infty(d,R)|I_-|
       \int_{G_\infty}\widehat K_\infty(y)\nu_\infty(\dd y)\\
&\leq B_\infty(d,R)|I_-|
       \mathbb E_\mu\widehat K_\infty(Y_\infty)\\
&\leq\bar\kappa_\infty B_\infty(d,R)|I_-|.
\end{aligned}
\]
An empty piece has an empty event, and a singleton has a singleton image of
zero length; finite-cap domination makes its conditional probability zero.
This is a separate negative-chart calculation and does not use a sign
symmetry.

Together, the three propositions in this subsection give the three matching
chart inequalities simultaneously.  Each outer piece uses the mean outer cap
once for its own event; no factor two has entered any individual estimate.
When $d=1$, Lemma~\ref{lem:kernel-caps} identifies the inner and outer labels
with the same unconditional coefficient law and the same finite density cap,
while $s_0=s_\infty=-\theta$ and $B_0(1,R)=B_\infty(1,R)=1$.  The inner
proposition nevertheless invokes the primitive bound with $\bar\kappa_0$,
and each outer proposition invokes the primitive bound with
$\bar\kappa_\infty$.  Thus coincidence of the degree-one kernel does not
merge the two cap hypotheses or either higher-degree conditioning interface.
\end{proof}

\subsection{Three-piece decomposition and exact maximum}
\label{app:three-piece}

\begin{lemma}[Three-piece parameter and event decomposition]
\label{lem:three-piece-decomposition}
Under Assumptions~\ref{assump:compact-parameter-domain} and
\ref{assump:indexed-regime}, for $I\in\mathcal I(\Theta)$ define
$I_0,I_+,I_-$ as in Section~\ref{sec:preliminaries}.  Then, with every
endpoint convention,
\[
I=I_0\mathbin{\dot\cup}I_+\mathbin{\dot\cup}I_-,
\qquad
|I|=|I_0|+|I_+|+|I_-|,
\]
the points $1$ and $-1$ belong only to $I_0$ when they belong to $I$, and
\[
H_{d,I}=H_{d,I_0}\cup H_{d,I_+}\cup H_{d,I_-}.
\]
\end{lemma}

\begin{proof}
The sets $[-1,1]$, $(1,\infty)$, and $(-\infty,-1)$ are pairwise disjoint
and partition $\R$.  Intersecting this exact partition with $I$ preserves
the endpoint convention and gives the disjoint union.  In particular, an
included $1$ or $-1$ occurs only in $I_0$; an excluded boundary point occurs
in none of the pieces.  Finite additivity of Lebesgue measure gives the
length identity, including empty and singleton pieces.

For a fixed coefficient vector, a root in $I$ lies in exactly one of the
three pieces, proving the forward event inclusion.  A root in any piece is a
root in $I$, proving the reverse inclusion.  The three events need not be
disjoint, since one polynomial can have roots in multiple pieces; the identity
is an event union, not a probability-additivity assertion.
\end{proof}

\begin{proposition}[Weighted chart union bound]
\label{prop:weighted-chart-bound}
Under Assumptions~\ref{assump:compact-parameter-domain} and
\ref{assump:indexed-regime}, Lemma~\ref{lem:three-piece-decomposition}, and
Propositions~\ref{prop:inner-chart-probability},
\ref{prop:positive-chart-probability}, and
\ref{prop:negative-chart-probability}, for every
$\mu\in\mathcal D_{d,R,\eta}$ and $I\in\mathcal I(\Theta)$,
\[
\mu(H_{d,I})
\leq\bar\kappa_0B_0(d,R)|I_0|
 +\bar\kappa_\infty B_\infty(d,R)(|I_+|+|I_-|).
\]
If a piece is empty or a singleton, its event probability and weighted term
are zero.
\end{proposition}

\begin{proof}
By Lemma~\ref{lem:three-piece-decomposition} and finite subadditivity,
\[
\mu(H_{d,I})
\leq\mu(H_{d,I_0})+\mu(H_{d,I_+})+\mu(H_{d,I_-}).
\]
Applying the three chart propositions gives
\[
\begin{aligned}
\mu(H_{d,I})
&\leq\bar\kappa_0B_0(d,R)|I_0|
 +\bar\kappa_\infty B_\infty(d,R)|I_+|\\
&\qquad+\bar\kappa_\infty B_\infty(d,R)|I_-|,
\end{aligned}
\]
which is the claimed factored form.  Empty and singleton cases are already
handled in the three preceding propositions and in the length identity; in
particular, singleton boundary pieces at $\pm1$ require no correction.
\end{proof}

\begin{proposition}[Exact two-chart maximum bound]
\label{prop:exact-chart-maximum}
Under Assumptions~\ref{assump:compact-parameter-domain} and
\ref{assump:indexed-regime}, Lemma~\ref{lem:three-piece-decomposition}, and
Proposition~\ref{prop:weighted-chart-bound}, every
$\mu\in\mathcal D_{d,R,\eta}$ and $I\in\mathcal I(\Theta)$ satisfy
\[
\mu(H_{d,I})\leq M_\eta(d,R)|I|,
\qquad
M_\eta(d,R)=\max\{\bar\kappa_0B_0(d,R),
\bar\kappa_\infty B_\infty(d,R)\}.
\]
\end{proposition}

\begin{proof}
For this scalar calculation set
\[
A:=\bar\kappa_0B_0(d,R),\qquad
C:=\bar\kappa_\infty B_\infty(d,R),\qquad
x:=|I_0|,\ y:=|I_+|,\ z:=|I_-|.
\]
The cap parameters are nonnegative and $d\geq1$, $R\geq1$, so
$A,C,x,y,z\geq0$.  With $m=\max\{A,C\}$,
\[
Ax+C(y+z)\leq mx+m(y+z)=\max\{A,C\}(x+y+z).
\]
Lemma~\ref{lem:three-piece-decomposition} gives
$x+y+z=|I|$, and Proposition~\ref{prop:weighted-chart-bound} gives the
left-hand side.  Since $\max\{A,C\}=M_\eta(d,R)$, the conclusion follows.
This is the exact maximum of the chart coefficients; replacing it by $A+C$
would be a strictly weaker estimate and is not needed.
\end{proof}

\subsection{Class supremum and fixed-parameter polynomial specialization}
\label{app:class-and-rate}

\begin{proposition}[Class-supremum closure]
\label{prop:class-supremum}
Under Assumption~\ref{assump:indexed-regime} and
Proposition~\ref{prop:exact-chart-maximum}, for every $d\geq1$ and $R\geq1$,
\[
C_{\mathcal D_{d,R,\eta}}\leq M_\eta(d,R).
\]
This includes the convention that $C_{\mathcal D_{d,R,\eta}}=0$ if either
$\mathcal D_{d,R,\eta}$ or $\mathcal I(\Theta)$ is empty.
\end{proposition}

\begin{proof}
The quantities $B_0$ and $B_\infty$ are nonnegative for the indexed regime,
so $M_\eta(d,R)\geq0$.  If both indexing sets are nonempty, then for every
admissible $\mu$ and every $I$ the preceding proposition gives
\[
\frac{\mu(H_{d,I})}{|I|}\leq M_\eta(d,R),
\]
because $|I|>0$.  Taking first the interval supremum and then the law
supremum preserves this inequality and yields the definition of
$C_{\mathcal D_{d,R,\eta}}$.  If either indexing set is empty, the stipulated
value of $C_{\mathcal D_{d,R,\eta}}$ is zero, which is still at most the
nonnegative $M_\eta(d,R)$.  No witness or lower cap threshold is used.
\end{proof}

\begin{lemma}[Inner coefficient domination]
\label{lem:inner-polynomial-domination}
Under Assumption~\ref{assump:indexed-regime}, for every $d\geq1$ and $R\geq1$,
\[
\bar\kappa_0B_0(d,R)\leq P_\eta(d,R).
\]
\end{lemma}

\begin{proof}
By definition, $0\leq\bar\kappa_0\leq\bar\kappa_*$.  Therefore
\[
\begin{aligned}
\bar\kappa_0B_0(d,R)
&=\bar\kappa_0d+\frac{\bar\kappa_0}{2}Rd(d-1)\\
&\leq\bar\kappa_*d+\frac{\bar\kappa_*}{2}Rd(d-1)\\
&\leq\bar\kappa_*d+\frac{\bar\kappa_*}{2}Rd^2
=P_\eta(d,R),
\end{aligned}
\]
where the second inequality is the pointwise comparison
$d(d-1)\leq d^2$.  All multipliers are nonnegative, including at
$\bar\kappa_*=0$.
\end{proof}

\begin{lemma}[Outer coefficient domination]
\label{lem:outer-polynomial-domination}
Under Assumption~\ref{assump:indexed-regime}, for every $d\geq1$ and $R\geq1$,
\[
\bar\kappa_\infty B_\infty(d,R)\leq P_\eta(d,R).
\]
\end{lemma}

\begin{proof}
Since $0\leq\bar\kappa_\infty\leq\bar\kappa_*$, $1\leq d$, and
$d(d-1)\leq d^2$, we have
\[
\begin{aligned}
\bar\kappa_\infty B_\infty(d,R)
&=\bar\kappa_\infty+\frac{\bar\kappa_\infty}{2}Rd(d-1)\\
&\leq\bar\kappa_*+\frac{\bar\kappa_*}{2}Rd(d-1)\\
&\leq\bar\kappa_*d+\frac{\bar\kappa_*}{2}Rd^2
=P_\eta(d,R).
\end{aligned}
\]
The inequality remains valid at the zero-cap boundary.
\end{proof}

\begin{proposition}[Fixed-$\eta$ polynomial specialization]
\label{prop:fixed-eta-polynomial-specialization}
Under Assumption~\ref{assump:indexed-regime} and
Lemmas~\ref{lem:inner-polynomial-domination} and
\ref{lem:outer-polynomial-domination}, for every $d\geq1$ and $R\geq1$,
\[
M_\eta(d,R)\leq P_\eta(d,R)
=\bar\kappa_*d+\frac{\bar\kappa_*}{2}Rd^2.
\]
\end{proposition}

\begin{proof}
The two preceding lemmas bound both entries of the maximum defining
$M_\eta(d,R)$ by the same $P_\eta(d,R)$.  Maximum monotonicity therefore
gives
\[
\max\{\bar\kappa_0B_0(d,R),\bar\kappa_\infty B_\infty(d,R)\}
\leq\max\{P_\eta(d,R),P_\eta(d,R)\}=P_\eta(d,R).
\]
This is pointwise for every allowed index, not asymptotic.  For fixed $\eta$,
$P_\eta$ has monomials $d$ and $Rd^2$, total degree at most three in
$(d,R)$, and its displayed coefficients depend only on $\eta$.  If
$\bar\kappa_*=0$, it is the zero polynomial and the same comparison holds.
The proof also verifies every technical condition used by the public
fixed-cap corollary: $d\geq1$, $R\geq1$, nonnegative cap parameters, and the
two explicit inequalities $1\leq d$ and $d(d-1)\leq d^2$.  The passage from
the per-law inequality to the class constant is deterministic and was
already proved in Proposition~\ref{prop:class-supremum}; no probability,
horizon, or norm conversion is hidden in this specialization.
\end{proof}

\subsection{Witness kernels, support, dependence, and membership}
\label{app:witness}

\begin{proposition}[Degree-one endpoint kernels]
\label{prop:degree-one-witness-kernels}
Under Assumption~\ref{assump:indexed-regime}, if $d=1$ and
$\mu^{\mathrm{wit}}_{1,R}$ is the witness law defined above, then both
regular conditional laws in the definition of
$\mathcal D_{1,R,\eta}$ are the marginal law
\[
\nu_R(A):=\frac{\lambda(A\cap[-R,R])}{2R}
=\int_Au_R(x)\dd x,
\qquad
u_R(x):=\frac1{2R}\1_{[-R,R]}(x).
\]
Consequently,
\[
K_0^{\mu^{\mathrm{wit}}_{1,R}}
=K_\infty^{\mu^{\mathrm{wit}}_{1,R}}
=\frac1{2R}
\quad\text{almost surely}.
\]
\end{proposition}

\begin{proof}
For $d=1$, the sole coefficient is uniform on $[-R,R]$, so its law and
density are exactly $\nu_R$ and $u_R$.  Both conditioning tuples are empty;
equivalently, the conditioning space is the one-point space $\R^0$, and the
conditional law at that point is the marginal $\nu_R$.  Since $R\geq1$,
$[-R,R]$ has positive length and
\[
\|u_R\|_{L^\infty(\R)}=\frac1{2R}.
\]
Every other density version agrees with $u_R$ almost everywhere and has the
same essential supremum.  Both setting-defined caps therefore equal the
displayed deterministic value.
\end{proof}

\begin{proposition}[Degree-two full-complement kernels]
\label{prop:degree-two-witness-kernels}
Under Assumption~\ref{assump:indexed-regime}, if $d=2$ and
$\mu^{\mathrm{wit}}_{2,R}$ is the witness law, then a regular conditional
law of $\alpha_0$ given $\alpha_1$ and a regular conditional law of
$\alpha_1$ given $\alpha_0$ are both the constant kernel $\nu_R$, whose
density is $u_R=(2R)^{-1}\1_{[-R,R]}$.  Consequently,
\[
K_0^{\mu^{\mathrm{wit}}_{2,R}}
=K_\infty^{\mu^{\mathrm{wit}}_{2,R}}
=\frac1{2R}
\quad\text{almost surely}.
\]
\end{proposition}

\begin{proof}
Here $\alpha_0$ and $\alpha_1$ are independent with common law $\nu_R$.
For every Borel $A\subseteq\R$, define the constant kernels
\[
Q^0_y(A):=\nu_R(A),\qquad Q^\infty_x(A):=\nu_R(A)
\quad(x,y\in\R).
\]
They are measurable in their conditioning variables.  For Borel
$A,C\subseteq\R$, independence gives
\[
\begin{aligned}
\mathbb P(\alpha_0\in A,\alpha_1\in C)
&=\nu_R(A)\mathbb P(\alpha_1\in C)
=\int_CQ^0_y(A)\mathbb P_{\alpha_1}(\dd y),\\
\mathbb P(\alpha_1\in A,\alpha_0\in C)
&=\nu_R(A)\mathbb P(\alpha_0\in C)
=\int_CQ^\infty_x(A)\mathbb P_{\alpha_0}(\dd x).
\end{aligned}
\]
Thus they are regular conditional laws after conditioning on the full
nonpivot coordinate.  Each has density $u_R$ and essential supremum
$1/(2R)$.  The integral identities determine regular conditional laws on a
full conditioning-marginal set; Radon--Nikodym uniqueness then makes the
density norm independent of the permissible versions.  Hence both random
caps have the asserted value almost surely.
\end{proof}

\begin{proposition}[Higher-degree full-complement kernels]
\label{prop:higher-degree-witness-kernels}
Under Assumption~\ref{assump:indexed-regime}, if $d\geq3$ and
$\mu^{\mathrm{wit}}_{d,R}$ is the witness law, then a regular conditional
law of $\alpha_0$ given all coordinates $\alpha_{1:d-1}$ and a regular
conditional law of $\alpha_{d-1}$ given all coordinates
$\alpha_{0:d-2}$ are both the constant uniform kernel $\nu_R$, with density
$u_R=(2R)^{-1}\1_{[-R,R]}$.  Consequently,
\[
K_0^{\mu^{\mathrm{wit}}_{d,R}}
=K_\infty^{\mu^{\mathrm{wit}}_{d,R}}
=\frac1{2R}
\quad\text{almost surely}.
\]
\end{proposition}

\begin{proof}
The construction is
\[
\alpha_0=U_0,\qquad
\alpha_j=RS\ (1\leq j\leq d-2),\qquad
\alpha_{d-1}=U_\infty,
\]
where $U_0,U_\infty\sim\nu_R$, $S$ is Rademacher, and all three variables
are mutually independent.  The complete conditioning tuples are
\[
\alpha_{1:d-1}=(RS,\ldots,RS,U_\infty),\qquad
\alpha_{0:d-2}=(U_0,RS,\ldots,RS).
\]
The first is a measurable function of $(S,U_\infty)$ and is independent of
$U_0=\alpha_0$.  The second is a measurable function of $(U_0,S)$ and is
independent of $U_\infty=\alpha_{d-1}$.  Although either complement reveals
$S$ on the witness support, it does not reveal the corresponding endpoint
pivot.

For all $y,z\in\R^{d-1}$, define
\[
Q^0_y(A):=\nu_R(A),\qquad Q^\infty_z(A):=\nu_R(A).
\]
For Borel $C\subseteq\R^{d-1}$, the two independence statements give
\[
\begin{aligned}
\mathbb P(\alpha_0\in A,\alpha_{1:d-1}\in C)
&=\nu_R(A)\mathbb P(\alpha_{1:d-1}\in C)\\
&=\int_CQ^0_y(A)\mathbb P_{\alpha_{1:d-1}}(\dd y),
\end{aligned}
\]
and
\[
\mathbb P(\alpha_{d-1}\in A,\alpha_{0:d-2}\in C)
=\int_CQ^\infty_z(A)\mathbb P_{\alpha_{0:d-2}}(\dd z).
\]
The constant kernels are measurable, so these identities verify them as
regular conditional laws on the full complement spaces, even at conditioning
values outside the witness support.  Their common density is $u_R$, whose
essential supremum is $1/(2R)$.  Almost-sure uniqueness of conditional laws
and of Radon--Nikodym densities shows that every permissible version has the
same realized cap.  Thus the two caps are exactly, rather than merely at
most, $1/(2R)$.
\end{proof}

\begin{proposition}[Cube support and ambient singularity]
\label{prop:witness-support-singularity}
Under Assumption~\ref{assump:indexed-regime}, the witness satisfies
\[
\mu^{\mathrm{wit}}_{d,R}([-R,R]^d)=1
\quad\text{for every }d\geq1.
\]
For every $d\geq3$, it is singular with respect to $d$-dimensional Lebesgue
measure.
\end{proposition}

\begin{proof}
For $d=1$, the only coordinate is uniform on $[-R,R]$.  For $d=2$, both
coordinates are uniform on this interval.  For $d\geq3$, the two endpoints
lie in $[-R,R]$ almost surely and every middle coordinate is
$RS\in\{-R,R\}$.  This proves cube support in all three regimes.

For $d\geq3$, set
\[
A_{d,R}:=[-R,R]\times
\{(R,\ldots,R),(-R,\ldots,-R)\}\times[-R,R]\subseteq\R^d,
\]
where each middle vector has $d-2\geq1$ entries.  This is a closed Borel set
and $\mu^{\mathrm{wit}}_{d,R}(A_{d,R})=1$.  Its middle factor is a finite
subset of $\R^{d-2}$ and has $(d-2)$-dimensional Lebesgue measure zero.
Product measure on the two product rectangles comprising $A_{d,R}$ gives
\[
\lambda_d(A_{d,R})=0.
\]
Thus a Lebesgue-null Borel set carries all witness mass.  The argument
includes $d=3$, where the middle factor is the two-point set $\{-R,R\}$.
\end{proof}

\begin{proposition}[Dependence of the middle coordinates]
\label{prop:witness-middle-dependence}
Under Assumption~\ref{assump:indexed-regime}, if $d\geq4$, then the middle
coordinates of $\mu^{\mathrm{wit}}_{d,R}$ are dependent; in particular,
$\alpha_1$ and $\alpha_2$ are not independent.
\end{proposition}

\begin{proof}
Both coordinates exist when $d\geq4$, and
$\alpha_1=RS=\alpha_2$ almost surely.  Since $R\geq1$, the values $R$ and
$-R$ are distinct.  Hence
\[
\mathbb P(\alpha_1=R)=\mathbb P(\alpha_2=R)=\frac12,
\qquad
\mathbb P(\alpha_1=R,\alpha_2=R)=\frac12.
\]
The joint probability is not the product $1/4$ of the marginals, proving
dependence.
\end{proof}

\begin{lemma}[Witness-cap threshold comparison]
\label{lem:witness-threshold}
Under Assumption~\ref{assump:indexed-regime}, suppose that the applicable
one of Propositions~\ref{prop:degree-one-witness-kernels},
\ref{prop:degree-two-witness-kernels}, and
\ref{prop:higher-degree-witness-kernels} gives
\[
K_0^{\mu^{\mathrm{wit}}_{d,R}}
=K_\infty^{\mu^{\mathrm{wit}}_{d,R}}
=\frac1{2R}\quad\text{almost surely}.
\]
If, only for this witness clause,
$\bar\kappa_0,\bar\kappa_\infty\geq1/2$, then
\[
\mathbb E K_0^{\mu^{\mathrm{wit}}_{d,R}}\leq\bar\kappa_0,
\qquad
\mathbb E K_\infty^{\mu^{\mathrm{wit}}_{d,R}}\leq\bar\kappa_\infty.
\]
\end{lemma}

\begin{proof}
Each cap is the deterministic value $1/(2R)$.  Since $R\geq1$,
\[
\mathbb E K_0^{\mu^{\mathrm{wit}}_{d,R}}
=\mathbb E K_\infty^{\mu^{\mathrm{wit}}_{d,R}}
=\frac1{2R}
\leq\frac12
\leq\min\{\bar\kappa_0,\bar\kappa_\infty\}.
\]
This is the only use of either lower threshold on the cap parameters.
\end{proof}

\begin{proposition}[Witness membership]
\label{prop:witness-membership}
Under Assumption~\ref{assump:indexed-regime} and
Propositions~\ref{prop:degree-one-witness-kernels}--\ref{prop:witness-middle-dependence}
and Lemma~\ref{lem:witness-threshold}, if
$\bar\kappa_0,\bar\kappa_\infty\geq1/2$, then for every $d\geq1$ and
$R\geq1$,
\[
\mu^{\mathrm{wit}}_{d,R}\in\mathcal D_{d,R,\eta},
\qquad
K_0^{\mu^{\mathrm{wit}}_{d,R}}
=K_\infty^{\mu^{\mathrm{wit}}_{d,R}}
=\frac1{2R}\quad\text{almost surely}.
\]
The witness is ambiently singular for $d\geq3$ and has dependent middle
coordinates for $d\geq4$.
\end{proposition}

\begin{proof}
Fix $(d,R)$ and impose the lower thresholds only for this membership
statement.

If $d=1$, Proposition~\ref{prop:degree-one-witness-kernels} supplies the two
conditional densities and exact caps for the empty conditioning tuples.
Proposition~\ref{prop:witness-support-singularity} supplies support on
$[-R,R]$, and Lemma~\ref{lem:witness-threshold} supplies both mean-cap
inequalities.

If $d=2$, Proposition~\ref{prop:degree-two-witness-kernels} supplies the two
full-complement conditional densities and caps, now conditioning each
endpoint on the other.  Proposition~\ref{prop:witness-support-singularity}
gives support on $[-R,R]^2$, and Lemma~\ref{lem:witness-threshold} supplies
the mean inequalities.

If $d\geq3$, Proposition~\ref{prop:higher-degree-witness-kernels} supplies
both full-complement conditional densities and caps even though the
complement contains the Rademacher middle block.
Proposition~\ref{prop:witness-support-singularity} gives cube support and
ambient singularity.  When $d\geq4$,
Proposition~\ref{prop:witness-middle-dependence} additionally proves the
middle-coordinate dependence.  Lemma~\ref{lem:witness-threshold} again gives
the two mean-cap inequalities.

In every degree regime, the witness is the law of a Borel measurable
Euclidean random vector and hence is a Borel probability law on $\R^d$.  The
preceding results verify every item in the definition of
$\mathcal D_{d,R,\eta}$: cube support, almost-sure existence of both endpoint
conditional densities, and both mean-cap inequalities.  This proves
membership and the exact cap identities.  Singularity and middle dependence
show that neither an ambient joint density nor coordinate independence has
been used as a hidden membership condition.  The threshold comparison is
isolated in Lemma~\ref{lem:witness-threshold} and has no role in the general
root-hitting argument.
\end{proof}

\subsection{Proof of the Main Theorem}
\label{app:main-proof}

\begin{proof}[Proof of Theorem~\ref{thm:main}]
Fix the finite $\eta$ first, and then fix arbitrary $d\geq1$, $R\geq1$,
$\mu\in\mathcal D_{d,R,\eta}$, and $I\in\mathcal I(\Theta)$.
Proposition~\ref{prop:weighted-chart-bound} gives the exact weighted
three-piece estimate
\[
\mu(H_{d,I})
\leq\bar\kappa_0B_0(d,R)|I_0|
 +\bar\kappa_\infty B_\infty(d,R)(|I_+|+|I_-|).
\]
Proposition~\ref{prop:exact-chart-maximum} then gives
\[
\mu(H_{d,I})\leq M_\eta(d,R)|I|.
\]
The first proposition uses only the finite union bound after the exact event
decomposition, and the second uses the scalar maximum inequality, so the
coefficient is the maximum of the two chart constants rather than their sum.

Proposition~\ref{prop:class-supremum} divides by the positive interval length
and takes the interval and law suprema, while also treating either empty
indexing set.  Hence
\[
C_{\mathcal D_{d,R,\eta}}\leq M_\eta(d,R).
\]
Proposition~\ref{prop:fixed-eta-polynomial-specialization} proves, pointwise
for every allowed $d,R$,
\[
M_\eta(d,R)\leq P_\eta(d,R)
=\bar\kappa_*d+\frac{\bar\kappa_*}{2}Rd^2.
\]
This completes the root-hitting and class-uniform conclusions without any
lower bound on $\bar\kappa_0$ or $\bar\kappa_\infty$.

Under the separate hypothesis
$\bar\kappa_0,\bar\kappa_\infty\geq1/2$,
Proposition~\ref{prop:witness-membership} gives, simultaneously for every
$d\geq1$ and $R\geq1$,
\[
\mu^{\mathrm{wit}}_{d,R}\in\mathcal D_{d,R,\eta},
\qquad
K_0^{\mu^{\mathrm{wit}}_{d,R}}
=K_\infty^{\mu^{\mathrm{wit}}_{d,R}}
=\frac1{2R}\quad\text{almost surely}.
\]
This proves the independent nonemptiness clause.

Every constant used above is one of the displayed setting-derived
$B_0,B_\infty,M_\eta,P_\eta$.  The chart estimates are per-law deterministic
probability inequalities, and the remaining operations are deterministic
finite addition, scalar domination, and suprema.  Thus no confidence
parameter, horizon upgrade, limiting argument, norm conversion, or hidden
dependence is introduced.  The exposed quantities and the static
Lebesgue-length and conditional-$L^\infty$ modes are exactly those stated in
Theorem~\ref{thm:main}.
\end{proof}
